% Options for packages loaded elsewhere
\PassOptionsToPackage{unicode}{hyperref}
\PassOptionsToPackage{hyphens}{url}
\PassOptionsToPackage{dvipsnames,svgnames,x11names}{xcolor}
\documentclass[
]{article}
\usepackage{xcolor}
\usepackage[margin = 0.5in]{geometry}
\usepackage{amsmath,amssymb}
\setcounter{secnumdepth}{-\maxdimen} % remove section numbering
\usepackage{iftex}
\ifPDFTeX
  \usepackage[T1]{fontenc}
  \usepackage[utf8]{inputenc}
  \usepackage{textcomp} % provide euro and other symbols
\else % if luatex or xetex
  \usepackage{unicode-math} % this also loads fontspec
  \defaultfontfeatures{Scale=MatchLowercase}
  \defaultfontfeatures[\rmfamily]{Ligatures=TeX,Scale=1}
\fi
\usepackage{lmodern}
\ifPDFTeX\else
  % xetex/luatex font selection
\fi
% Use upquote if available, for straight quotes in verbatim environments
\IfFileExists{upquote.sty}{\usepackage{upquote}}{}
\IfFileExists{microtype.sty}{% use microtype if available
  \usepackage[]{microtype}
  \UseMicrotypeSet[protrusion]{basicmath} % disable protrusion for tt fonts
}{}
\makeatletter
\@ifundefined{KOMAClassName}{% if non-KOMA class
  \IfFileExists{parskip.sty}{%
    \usepackage{parskip}
  }{% else
    \setlength{\parindent}{0pt}
    \setlength{\parskip}{6pt plus 2pt minus 1pt}}
}{% if KOMA class
  \KOMAoptions{parskip=half}}
\makeatother
\usepackage{graphicx}
\makeatletter
\newsavebox\pandoc@box
\newcommand*\pandocbounded[1]{% scales image to fit in text height/width
  \sbox\pandoc@box{#1}%
  \Gscale@div\@tempa{\textheight}{\dimexpr\ht\pandoc@box+\dp\pandoc@box\relax}%
  \Gscale@div\@tempb{\linewidth}{\wd\pandoc@box}%
  \ifdim\@tempb\p@<\@tempa\p@\let\@tempa\@tempb\fi% select the smaller of both
  \ifdim\@tempa\p@<\p@\scalebox{\@tempa}{\usebox\pandoc@box}%
  \else\usebox{\pandoc@box}%
  \fi%
}
% Set default figure placement to htbp
\def\fps@figure{htbp}
\makeatother
\setlength{\emergencystretch}{3em} % prevent overfull lines
\providecommand{\tightlist}{%
  \setlength{\itemsep}{0pt}\setlength{\parskip}{0pt}}
\usepackage{color}
\usepackage{framed}
\setlength{\fboxsep}{.8em}

\usepackage{tcolorbox}

\tcbset{
  mystyle/.style={
    before skip=-10pt, % Adjust this value
    after skip=-10pt   % Adjust this value
  }
}
\newtcolorbox{mynoteblock}[1][]{mystyle, #1}

\newtcolorbox{blackbox}{
  colback=black,
  colframe=orange,
  coltext=white,
  boxsep=5pt,
  arc=4pt}
  
\newtcolorbox{ltbluebox}{
  colback=blue!5!white,
  colframe=blue!75!black,
  boxsep=5pt,
  arc=4pt}
  
\newtcolorbox{ltgreenbox}{
  colback=green!5!white,
  colframe=blue!75!black,
  boxsep=5pt,
  arc=4pt}
  
\newtcolorbox{ltpurplebox}{
  colback=purple!5!white,
  colframe=blue!75!black,
  boxsep=5pt,
  arc=4pt}
  


\newenvironment{cols}[1][]{}{}

\newenvironment{col}[1]{\begin{minipage}[t]{#1}\ignorespaces\vspace{0pt}}{%
\end{minipage}
\ifhmode\unskip\fi
\aftergroup\useignorespacesandallpars}

\def\useignorespacesandallpars#1\ignorespaces\fi{%
#1\fi\ignorespacesandallpars}

\makeatletter
\def\ignorespacesandallpars{%
  \@ifnextchar\par
    {\expandafter\ignorespacesandallpars\@gobble}%
    {}%
}
\makeatother

\usepackage{awesomebox}
\setlength{\aweboxvskip}{1mm}


\pagenumbering{gobble}

\usepackage{enumitem}
\setlist{itemsep=1pt, parsep=0pt} %# Adjust vertical spacing
\setlist[itemize,1]{leftmargin=*, labelsep=0.3em} %# Adjust horizontal spacing
\setlist[enumerate,1]{leftmargin=*, labelsep=0.3em} %# Adjust horizontal spacing

\usepackage{bookmark}
\IfFileExists{xurl.sty}{\usepackage{xurl}}{} % add URL line breaks if available
\urlstyle{same}
\hypersetup{
  colorlinks=true,
  linkcolor={Maroon},
  filecolor={Maroon},
  citecolor={Blue},
  urlcolor={blue},
  pdfcreator={LaTeX via pandoc}}

\author{}
\date{\vspace{-2.5em}}

\begin{document}

\vspace{-1cm}

\section{Measures Setting: Potential Indicators and
Pathways}\label{measures-setting-potential-indicators-and-pathways}

\begin{noteblock}
Measure setting is the process of establishing regulatory controls and
conservation tools that govern how, when, where, and how much fish can
be harvested to prevent overfishing and ensure the long-term
sustainability of stocks. \textbf{Here we emphasize measures governing
how, when, and where to harvest, because how much to harvest is
described separately under Catch/Quota setting.} Measures could include
regulations such as gear rewquirements, spatial or temporal openings and
closures, and size and age limits.

\end{noteblock}

\begin{itemize}
\tightlist
\item
  \textbf{Key question: What indicators could inform measures setting
  across all stocks? Could ecosystem indicators or recreational demand
  indicators in addition to current information improve the process?}
\item
  \textbf{Scope: Should we develop a general process evaluating
  multifleet objectives in the measures setting process, or focus on an
  application for a single species? Which species?}
\end{itemize}

\begin{cols}

\begin{col}{0.49\textwidth}

\begin{ltbluebox}

\subsection{What information is used
now?}\label{what-information-is-used-now}

Current measures in place are outlined by NOAA's
\href{https://www.fisheries.noaa.gov/species-directory}{Specification
Summaries} and their
\href{https://www.fisheries.noaa.gov/rules-and-regulations/fisheries}{Federal
Regulations}.

\vspace{2ex}

Measures are set for individual stocks within FMPs to ``establish
allowable harvest levels that will prevent overfishing, consistent with
the most recent scientific information.''

\begin{itemize}
\tightlist
\item
  Commercial measures include gear requirements, seasons, and areas. The
  goals of measures are listed in FMPs, for example Summer Flounder FMP
  Objective 3.1 is to ``Provide reasonable access to the fishery
  throughout the management unit. Fishery allocations and other
  management measures should balance responsiveness to changing social,
  economic, and ecological conditions with historic and current
  importance to various user groups and communities.''
\item
  Recreational measures include recreational bag, size, and season
  limits. Recreational measures are designed to achieve overall harvest
  within the Recreational Harvest Limit (RHL) as well as other goals
  related to angler welfare.
\item
  Measures can be changed by either FMP amendment or framework action.
\end{itemize}

\vspace{2ex}

Workshop 1 discussions centered on maximizing fishing opportunities
across species and fleets when setting measures for individual species.
\textbf{Opportunities across species are not typically considered when
setting measures for individual species.}

\begin{ltpurplebox}
This emphasizes the multifleet/multispecies aspect highlighted in the
workshop to ensure overall opportunities are maximized, but ecosystem
indicators could also be used in single species measures without
considering this (e.g.~projected conditions impact on growth could lead
to different fish size or gear recommendations, changing seasonality of
migration or spawning could lead to different open/closed seasons, etc)
simpler but less focused on a system level objective.

\end{ltpurplebox}

\end{ltbluebox}

\end{col}

\begin{col}{0.02\textwidth}
~

\end{col}

\begin{col}{0.49\textwidth}

\begin{ltgreenbox}

\subsection{Information sources for
indicators}\label{information-sources-for-indicators}

To understand multispecies, multifleet implications of measures set for
a given species, we need:

\begin{itemize}
\item
  What fleets the species is caught in.

  \begin{itemize}
  \tightlist
  \item
    \href{https://www.fisheries.noaa.gov/new-england-mid-atlantic/commercial-fishing/quota-monitoring-greater-atlantic-region}{CAMS}
    data for commercial fleets
  \item
    \href{https://www.fisheries.noaa.gov/recreational-fishing-data/introduction-marine-recreational-information-program-data}{MRIP}
    for recreational fleets
  \end{itemize}
\item
  Other targets and catch from these fleets
\item
  Who manages the other species and fleets if not MAFMC
\item
  Measures in place for all the other species (commercial and
  recreational) caught with the species of interest
\end{itemize}

\vspace{2ex}

Big picture performance metrics for commercial and recreational
fisheries (landings, revenues, trips, and other system level fishery
attributes) are available in:

\begin{itemize}
\item
  \href{https://www.fisheries.noaa.gov/new-england-mid-atlantic/ecosystems/state-ecosystem-reports-northeast-us-shelf}{State
  of the Ecosystem Reports}
\item
  \href{https://static1.squarespace.com/static/511cdc7fe4b00307a2628ac6/t/672bbef75e7b1166ec15dba9/1730920200863/2_MAB_RiskAssess_2024.pdf}{EAFM
  Risk Assessments}
\item
  \href{https://www.fisheries.noaa.gov/new-england-mid-atlantic/ecosystems/ecosystem-and-socioeconomic-profile-development-and-reports}{Ecosystem
  and Socioeconomic Profiles}
\item
  Stock assessment projections of age and size compositions
\item
  Rec Demand Model projections of commercial fleet and angler response
  to measures

  \begin{itemize}
  \tightlist
  \item
    Threshold for switching target species?
  \end{itemize}
\item
  Projections of climate and habitat conditions for relevant species

  \begin{itemize}
  \tightlist
  \item
    \href{https://psl.noaa.gov/cefi_portal/}{Changing Ecosystems and
    Fisheries Initiative}
  \end{itemize}
\item
  Projections of species overlap; see
  \href{https://www.rhapcast.org/}{river herring} bycatch mitigation IRA
  project
\end{itemize}

\subsubsection{Potential Pathways}\label{potential-pathways}

Develop a live view of all current measures by species and fleet that
shows effects of changes for one stock.

\begin{ltpurplebox}

\subsection{Information gaps}\label{information-gaps}

Aggregate performance metrics for the fisheries to evaluate implications
of changes (commercial and recreational):

\begin{itemize}
\tightlist
\item
  Expenditures
\item
  Resilience
\end{itemize}

Easily accessible, up to date measures in place for all species across
all states

\end{ltpurplebox}

\end{ltgreenbox}

\end{col}

\end{cols}

\end{document}
